\section{Introduction}\label{sec:intro}

Unmanned surface vehicles (USVs) carry out scientific, commercial, and defense missions in oceans, coastal waters, rivers, and lakes [liu2016usv]. USVs monitor the water quality, survey the bottom topography, patrol maritime areas, collect water samples, and transport cargo [liu2016usv, bai2022review, Peng2023]. A mission that covers a large area or that measures several points at the same time uses a fleet of USVs, and a ground station (GS) coordinates the vessels of the fleet [kim2026Command, couceiro2026Paradigm].

A digital twin (DT) is a digital representation of a physical object that exchanges data with this object automatically in both directions [kritzinger2018digital]. The DT of a fleet holds one twin per vehicle and updates each twin with the data of its vehicle [yan2026Swarm, zhang2024Formation], so a GS that runs the twins needs the state of every vessel in every period of the twins. The vessels and the GS exchange these states and the decisions of the twins over one wireless network, and the vessels contend for the channel of this network [Shang2025, bianchi2000performance]. A network digital twin (NDT) is a DT whose physical object is a data network [zhou2026ndt, almasan2022network], and an NDT evaluates a configuration of the network before the network applies it [almasan2022network].

The data of one vessel differ in size by orders of magnitude, from the state of the vessel, a few tens of bytes per update, to the frames of a stereo camera, tens of kilobytes per frame [zed]. When the telemetry, the states, and the video of a vessel share one Transmission Control Protocol (TCP) connection, which delivers the bytes in order [rfc9293], and the vessels offer more bytes than the channel carries, the states wait behind the video, and the GS receives the states late or in bursts [icecs2026fleetdt]. The twins of the GS then compute their decisions from late states.

The related works reduce the load of each vehicle at the source [mtowe2025Edge, xu2025Dynamic, kremer2023Teleoperated], reserve bandwidth for each MQTT flow in a wired network [shahri2022Extending], or model wired networks [almasan2022network, rusek2020routenet]. The DBA of passive optical networks and of IEEE~802.16, the controlled access of IEEE~802.11, and the overlay TDMA schemes over commodity IEEE~802.11 divide a shared medium among the stations in each cycle or slot [kramer2002ipact, ieee80216std2017, ieee80211std2020, rao2005overlay, djukic2009soft]. These works leave open the allocation of one shared Wi-Fi channel among the vessels of a fleet in each period, in the order of the data that the GS needs, with a bound on the fraction of the periods in which the GS misses an allocated message for any behavior of the channel. This paper addresses this allocation with three research questions (RQs):

\noindent RQ.I: under which condition on the capacity of the channel does the GS receive, within a period, every message that the allocation admits, and which fraction of the periods can the allocation guarantee for any behavior of the channel?

\noindent RQ.II: how accurately does a prediction from the busy time of the channel and a model of the IEEE~802.11 DCF estimate the capacity of the next period, compared with predictions from the measured capacity alone?

\noindent RQ.III: does the admission of the data of the fleet before each period keep the delivery of the data that the GS needs when the load of the fleet exceeds the capacity of the channel, compared with the prioritization of the IEEE~802.11 MAC, strict priority, and overlay TDMA?

\noindent Section~\ref{subsec:admission} answers RQ.I with Lemma~\ref{lem:delivery} and Theorem~\ref{thm:aci}, and the experiments of Section~\ref{sec:method} test RQ.I and answer RQ.II and RQ.III.

This paper presents Fleet-DT, an NDT that allocates one Wi-Fi channel among the vessels of a fleet in steps of $T_{\mathrm{s}} = 125$~ms by dynamic bandwidth allocation (DBA) above the IEEE~802.11 MAC. In each step, each vessel reports the demand of its data and the bytes that its MQTT bridge still holds; the GS predicts the capacity of the channel from the busy time that the access point measures and from the efficiency of the DCF, corrects the prediction with adaptive conformal inference (ACI) [gibbs2021adaptive], and admits the data of the step in the order of need of its applications within the corrected capacity; and the regulator in the MQTT broker of each vessel releases the admitted data and withholds the other data.

This paper extends the conference paper [icecs2026fleetdt]. With respect to [icecs2026fleetdt], this paper
(i) twins the network of the fleet with the DT prototype (DTP), the DT instance (DTI), the DT aggregate (DTA), and the DT environment (DTE) of Grieves and Vickers [grieves2017], with the state of the network of each step in the DTI and a what-if evaluation of configurations in the DTE (Section~\ref{subsec:dts});
(ii) defines the capacity of a step as a quantity that the access point measures and predicts it with the DCF model of Bianchi and the busy time of the channel (Section~\ref{subsec:pred});
(iii) admits the units of data of every vessel within one airtime budget in the order of need of the GS, with a required set of topics in every step and whole video frames, and proves that the admission returns the lexicographically largest feasible set (Section~\ref{subsec:admission});
(iv) proves that the GS receives every admitted message within the step when the work of the step fits in the realized capacity, and that ACI bounds the long-run fraction of the steps with a missed admitted message for every realization of the channel (Section~\ref{subsec:admission});
(v) states the requirements of a platform and the regulator that releases the admitted data at the pace of the radio (Section~\ref{sec:platform}); and
(vi) evaluates Fleet-DT with the unchanged software on emulated hosts over a real-time ns-3 emulation of the channel, with hypotheses, thresholds, and seeds fixed in a tagged commit before the evaluation (Sections~\ref{sec:method} and~\ref{sec:results}).

Section~\ref{sec:background} presents the concepts of the literature on which Fleet-DT builds and compares Fleet-DT with the related works, Section~\ref{sec:problem} defines the network and states the allocation problem, Section~\ref{sec:dt} defines the model of Fleet-DT, and Section~\ref{sec:platform} states the requirements of a platform. Section~\ref{sec:method} presents the experimental methodology, Section~\ref{sec:results} the results, and Section~\ref{sec:discussion} the discussion. Section~\ref{sec:conclusion} concludes the paper.
